\begin{figure}[H]
    \centering
    \begin{tikzpicture}[x=1.0cm, y=1.0cm]

        \def\bh{0.48cm}

        \node[rowlab, font=\scriptsize] at (-0.08,0.85) {\tl{host}: κωδικοποίηση};
        \node[rowlab, font=\scriptsize] at (-0.08,0.00) {\tl{NPU}: 8 πυρήνες \tl{AIE}};
        \node[rowlab, font=\scriptsize] at (-0.08,-0.85) {\tl{host}: συγχώνευση};

        \foreach \b/\lab in {0/1,1/2,2/3} {
            \pgfmathsetmacro{\xs}{\b*2.30}
            \node[blkdrm, minimum width=2.30cm, minimum height=\bh, anchor=west]
                at (\xs,0.85) {\scriptsize $b_{\lab}$};
            \pgfmathsetmacro{\xc}{\xs+2.30}
            \node[blknpu, minimum width=1.25cm, minimum height=\bh, anchor=west]
                at (\xc,0.00) {\scriptsize $b_{\lab}$};
            \pgfmathsetmacro{\xm}{\xs+4.60}
            \node[blknpu, minimum width=0.90cm, minimum height=\bh, anchor=west]
                at (\xm,-0.85) {\scriptsize $b_{\lab}$};
        }

        \node[notemute, font=\tiny] at (4.08,0.00) {αδράνεια};
        \node[notemute, font=\tiny] at (6.20,-0.85) {αδράνεια};

        \draw[dashed, draw=figmute, line width=0.4pt] (2.30,1.22) -- (2.30,-1.22);
        \draw[dashed, draw=figmute, line width=0.4pt] (4.60,1.22) -- (4.60,-1.22);

        \draw[bracemirror] (0,-1.22) -- (2.30,-1.22);
        \node[note, font=\scriptsize, anchor=north] at (1.15,-1.38) {περίοδος = πλατύτερο στάδιο};

        \draw[flow] (0,-2.05) -- (10.30,-2.05);
        \node[notemute, font=\tiny, anchor=north east] at (10.30,-2.12) {χρόνος};

    \end{tikzpicture}
    \caption{Η διοχέτευση τριών σταδίων της εκφόρτωσης. Ενώ τα πλακίδια
    επεξεργάζονται την παρτίδα $b_i$, ο \tl{host} κωδικοποιεί ήδη την $b_{i+1}$ στον
    δεύτερο \tl{buffer} και συγχωνεύει τα αποτελέσματα της $b_{i-1}$. Επειδή τα τρία
    στάδια επικαλύπτονται, ο συνολικός χρόνος δεν είναι το άθροισμά τους αλλά το
    γινόμενο του πλήθους των παρτίδων επί τη διάρκεια του πλατύτερου σταδίου. Στη
    συγκεκριμένη υλοποίηση το στάδιο αυτό είναι η κωδικοποίηση στον \tl{host}: μόνο
    αυτό εκτελείται χωρίς διακοπή, ενώ τα υπόλοιπα δύο μένουν αδρανή για μέρος κάθε
    περιόδου. Καμία περαιτέρω επιτάχυνση του πυρήνα \tl{AIE} δεν μπορεί επομένως να
    μειώσει τον συνολικό χρόνο.}
    \label{fig:npu-pipeline}
\end{figure}
